\subsection{Frozen Evaluation Results}
On the complete chronological HAI test split (4,443 prepared windows), the final
process-graph model achieved a ROC--AUC of 0.758, average precision of
0.376, attack-episode recall of
0.538, and eTaPR-F1 of 0.457. Relative to
the base detector, ROC--AUC increased from 0.688 to
0.758, while attack-episode recall increased from
0.231 to 0.538. These gains
were accompanied by a higher false-alarm window rate
(0.0047 versus 0.0121) and
worse expected calibration error (0.167 versus 0.246).

The primary frozen external red-team evaluation comprised 80 held-out tasks
and 320 Qwen3-1.7B candidates. After excluding one
fact-changing response and one no-op, 318
candidates were eligible. The legacy structural verifier accepted
177 candidates
(55.66\%) and exhibited a best-of-four task
attack-success rate of 78.75\%. The canonical
hardened verifier accepted no eligible candidate, with 95\% Wilson upper bounds of
1.19\% at candidate level and
4.58\% at task level. Clean
rejection remained zero. The paired exact McNemar test was significant
($p=2.168e-19$).

These results are bounded to the frozen HAI/Qwen3-1.7B protocol. Canonical rendering is a
secure-by-construction constraint on operator-facing wording rather than a universal
semantic detector. Moreover, automatic eligibility verifies committed facts and a
material surface change but does not establish harmful decision influence; blinded
human/domain-expert review remains necessary. eTaPR is reported at the prepared window
resolution and is not directly comparable with raw-second evaluations.
